% !TeX root = main.tex
% 第二章 质点动力学习题解答 子文件
% solution 环境在主文档main.tex中定义，本子文件不再重复定义
\section{质点动力学习题解答}
\subsection{喷气式飞机引擎推力}
某喷气式飞机以 $200\ \mathrm{m\cdot s^{-1}}$ 的速度在空中飞行，引擎每秒吸入 $50\ \mathrm{kg}$ 的空气与 $2\ \mathrm{kg\cdot s^{-1}}$ 的燃料混合燃烧，燃烧后的气体相对飞机以 $400\ \mathrm{m\cdot s^{-1}}$ 的速度向后喷出，试求该喷气式飞机引擎的推力。
\begin{solution}
取匀速飞行的飞机为惯性参考系，取 $\Delta t=1\ \mathrm{s}$。
$\Delta t$ 内吸入空气质量 $m_a=50\ \mathrm{kg}$，消耗燃料质量 $m_f=2\ \mathrm{kg}$。
空气进入引擎前相对飞机向前：$v_1=200\ \mathrm{m\cdot s^{-1}}$；
燃烧后气体相对飞机向后喷出：$v_2=-400\ \mathrm{m\cdot s^{-1}}$。
对气体应用动量定理，引擎对气体的冲量等于气体动量变化：
\[
F \Delta t = (m_a+m_f)v_2 - m_a v_1
\]
代入 $\Delta t=1\ \mathrm{s}$：
\[
F = (50+2)\cdot(-400) - 50 \cdot 200 = -30800\ \mathrm{N}
\]
$F$ 为引擎对气体的作用力，负号表示方向向后；根据牛顿第三定律，气体对引擎的推力向前：
\[
\boxed{F_\text{推}=30800\ \mathrm{N}=3.08\times 10^4\ \mathrm{N}}
\]
\end{solution}
\subsection{挖去方块平板质心（负质量法）}
一块边长 $L=6\ \mathrm{m}$ 的均匀方形平板，其一边被剪去边长 $d=2\ \mathrm{m}$ 的一块。被剪去一块的中心位于 $x=2\ \mathrm{m},\ y=0$ 处，完整平板中心在 $x=0,\ y=0$ 处，求平板剩余部分的质心位置。
\begin{solution}
设面密度为 $\sigma$。
完整大板质量：$M=\sigma L^2$；挖去小块质量：$m=\sigma d^2$。
负质量法：剩余物体 = 完整大板 + 负质量小块
\[
x_c = \frac{M\cdot x_M - m\cdot x_m}{M-m},\quad y_c = \frac{M\cdot y_M - m\cdot y_m}{M-m}
\]
已知 $x_M=0,\ y_M=0;\ x_m=2,\ y_m=0$
\[
x_c=\frac{\sigma L^2 \cdot 0 - \sigma d^2 \cdot 2}{\sigma(L^2-d^2)}
=\frac{-2d^2}{L^2-d^2}
\]
代入 $L=6,\ d=2$：
\[
x_c=\frac{-2\times 4}{36-4}=\frac{-8}{32}=-0.25\ \mathrm{m},\quad y_c=0
\]
\[
\boxed{(x_c,\,y_c)=(-0.25\ \mathrm{m},\ 0)}
\]
\end{solution}
\subsection{蓄水池抽水做功与抽水时间}
一长方体蓄水池，底面积 $S=50\ \mathrm{m^2}$，贮水深度 $h_1=1.5\ \mathrm{m}$。假定水平面距离地面的高度为 $h_2=5\ \mathrm{m}$。问要将这池水全部抽到地面上来，抽水机需做功多少？若抽水机的效率为 $80\%$，输入功率为 $P=35\ \mathrm{kW}$，则抽完这池水需要多长时间？
\begin{solution}
水总质量：$m=\rho S h_1$，水的重心初始高度：$h_c=h_2-\dfrac{h_1}{2}$。
抽水机有用功等于水重力势能增量：
\[
W_\text{有}=mg h_c=\rho S h_1 g \left(h_2-\frac{h_1}{2}\right)
\]
代入 $\rho=1000\ \mathrm{kg\cdot m^{-3}},\ g=9.8\ \mathrm{m\cdot s^{-2}},\ S=50\ \mathrm{m^2},\ h_1=1.5\ \mathrm{m},\ h_2=5\ \mathrm{m}$
\[
W_\text{有}=1000\times 50 \times 1.5 \times 9.8 \times \left(5-\frac{1.5}{2}\right)=3.18375\times 10^6\ \mathrm{J}
\]
效率定义 $\eta=\dfrac{W_\text{有}}{W_\text{总}}$，总功 $W_\text{总}=Pt$，故
\[
t=\frac{W_\text{有}}{\eta P}
\]
代入 $\eta=0.8,\ P=35\times10^3\ \mathrm{W}$
\[
t=\frac{3.18375\times 10^6}{0.8\times 35\times 10^3}\approx 113.7\ \mathrm{s}
\]
\[
\boxed{W_\text{有}=3.18\times 10^6\ \mathrm{J},\quad t\approx 114\ \mathrm{s}}
\]
\end{solution}
\subsection{变力拉力做功（积分法）}
用一变力 $F_\mathrm{T}=10\sin\alpha\ (\mathrm{N})$ 通过轻绳和轻滑轮将一物体从 A 处（$\alpha_1=30^\circ$）拉到 B 处（$\alpha_2=60^\circ$）。设 $h=2\ \mathrm{m}$，用积分法求拉力 $F_\mathrm{T}$ 在此过程中对物体所做的功。
\begin{solution}
几何关系：绳长 $l=\dfrac{h}{\sin\alpha}$，物体移动对应绳长缩短
\[
dl = -\frac{h \cos\alpha}{\sin^2\alpha} d\alpha
\]
拉力元功 $dW=F_T(-dl)$（$dl$ 为绳长减少，取绳缩短为正）：
\[
dW=F_T (-dl)=10\sin\alpha \cdot \frac{h \cos\alpha}{\sin^2\alpha} d\alpha=\frac{10h \cos\alpha}{\sin\alpha}d\alpha
\]
积分区间 $\alpha\colon 30^\circ \to 60^\circ$，即 $\pi/6 \to \pi/3$：
\[
W=\int_{\pi/6}^{\pi/3}\frac{10h \cos\alpha}{\sin\alpha}d\alpha
=10h \int_{\pi/6}^{\pi/3}\frac{d(\sin\alpha)}{\sin\alpha}
=10h \ln\sin\alpha \bigg|_{\pi/6}^{\pi/3}
\]
代入 $h=2\ \mathrm{m}$：
\[
W=20\left(\ln\frac{\sqrt{3}}{2}-\ln\frac12\right)=20\ln\sqrt{3}=10\ln 3 \approx 10.99\ \mathrm{J}
\]
\[
\boxed{W=10\ln 3 \approx 11.0\ \mathrm{J}}
\]
\end{solution}
\subsection{蒸汽打桩机（变阻力+完全非弹性碰撞）}
铁塔高 $40\ \mathrm{m}$，锤的质量 $M=10\ \mathrm{t}$。现将长达 $38.5\ \mathrm{m}$ 的钢筋混凝土桩打入地层，已知桩的质量为 $m=24\ \mathrm{t}$，横截面积为 $0.25\ \mathrm{m^2}$ 的正方形，桩的侧面单位面积所受泥土的阻力为 $F_r=2.65\times 10^4\ \mathrm{N\cdot m^{-2}}$。
\begin{enumerate}
    \item 靠桩自重能下沉多深？
    \item 桩稳定后，把锤提高 $1\ \mathrm{m}$，然后让锤自由下落而击桩，假定锤与桩发生完全非弹性碰撞，一锤能打下多深？
\end{enumerate}
\begin{solution}
桩截面正方形边长 $a=\sqrt{0.25}=0.5\ \mathrm{m}$，周长 $C=4a=2\ \mathrm{m}$；下沉深度 $s$ 时侧面积 $A=C s$，阻力 $f(s)=F_r \cdot C s$。
\paragraph{(1) 自重下沉深度 $s_0$}
动能定理，初末动能均为0：重力做功等于克服变阻力做功
\[
mg s_0 - \int_{0}^{s_0} F_r C s \,ds = 0
\]
\[
mg s_0 = \frac12 F_r C s_0^2 \implies s_0=\frac{2mg}{F_r C}
\]
代入 $m=24\times10^3\ \mathrm{kg},\ g=9.8\ \mathrm{m\cdot s^{-2}},\ F_r=2.65\times10^4\ \mathrm{N\cdot m^{-2}},\ C=2\ \mathrm{m}$
\[
s_0=\frac{2\times 24\times10^3\times9.8}{2.65\times10^4 \times 2}\approx 8.88\ \mathrm{m}
\]
\paragraph{(2) 锤下落撞击，一次打入深度 $\Delta s$}
锤自由下落 $h=1\ \mathrm{m}$：$v_0=\sqrt{2gh}$。
完全非弹性碰撞动量守恒：$M v_0=(M+m) v_\text{共}$
\[
v_\text{共}=\frac{M}{M+m}\sqrt{2gh}
\]
碰撞后下沉 $\Delta s$，动能定理：
\[
(M+m)g\Delta s - \int_{s_0}^{s_0+\Delta s} F_r C s\,ds = 0-\frac12(M+m)v_\text{共}^2
\]
\[
(M+m)g\Delta s - \frac12 F_r C \big[(s_0+\Delta s)^2-s_0^2\big] = -\frac{M^2 g h}{M+m}
\]
展开 $(s_0+\Delta s)^2-s_0^2=2s_0\Delta s+(\Delta s)^2$，代入数值得到关于 $\Delta s$ 的一元二次方程，取物理正根：
\[
\boxed{s_0\approx 8.88\ \mathrm{m},\quad \Delta s\approx 0.20\ \mathrm{m}}
\]
\end{solution}
\subsection{旋转匀质绳的内部张力}
一条质量为 $m$ 且分布均匀的绳子，长度为 $L$，一端拴在转轴上，并以恒定角速度 $\omega$ 在水平面上旋转，设转动过程中绳子始终伸直，且忽略重力与空气阻力，求距转轴为 $r$（$r<L$）处绳中的张力。
\begin{solution}
线密度 $\lambda = \dfrac{m}{L}$。
取距离转轴 $r$ 到 $L$ 的一段绳元作为研究对象。在位置 $x$ 处取微元 $dx$，微元质量 $dm=\lambda dx$，微元做圆周运动向心力 $dT = dm \cdot \omega^2 x$。
$r$ 处张力提供外侧全部绳段的向心力：
\[
T(r) = \int_{r}^{L} \lambda \omega^2 x \,dx
\]
\[
T(r) = \lambda \omega^2 \cdot \frac{x^2}{2}\bigg|_{r}^{L}
= \frac{m}{2L}\omega^2\left(L^2-r^2\right)
\]
\[
\boxed{T(r) = \frac{m\omega^2}{2L}\left(L^2-r^2\right)}
\]
\end{solution}
\subsection{匀质细棒浸入液体下落（动能定理）}
在密度为 $\rho_1$ 的液体上方悬一长为 $l$，密度为 $\rho_2$ 的均匀细棒，棒的 $B$ 端刚好和液面接触。今剪断细绳，设棒只在浮力和重力作用下下沉，
\begin{enumerate}
  \item 棒刚好全部浸入液体时的速度；
  \item 若 $\rho_2<\rho_1/2$ 时，求棒浸入液体的最大深度；
  \item 棒下落过程中所能达到的最大速率。
\end{enumerate}
\begin{solution}
设棒横截面积为 $S$，总质量 $m=\rho_2 S l$。
\begin{enumerate}
\item \textbf{刚好全部浸入时速度（下落距离 $l$）}
重力做功：$W_G = mg \cdot \dfrac{l}{2}$（重心下移 $l/2$）。
浸入深度 $x$ 时浮力 $F_\text{浮}(x)=\rho_1 g S x$，浮力做负功：
\[
W_\text{浮} = -\int_{0}^{l} \rho_1 g S x \,dx = -\frac12 \rho_1 g S l^2
\]
动能定理：$W_G + W_\text{浮} = \dfrac12 m v_1^2 - 0$
\[
\rho_2 S l g \cdot \frac{l}{2} - \frac12 \rho_1 g S l^2
= \frac12 \rho_2 S l \,v_1^2
\]
约去公共项整理：
\[
\rho_2 g l^2 - \rho_1 g l^2 = \rho_2 l\, v_1^2
\Rightarrow
v_1=\sqrt{\frac{(\rho_2-\rho_1)gl^2}{\rho_2 l}}
=\sqrt{\frac{(\rho_2-\rho_1)gl}{\rho_2}}
\]
注意：若 $\rho_2<\rho_1$，根号内为负，物理含义：棒不可能完全浸入；本题模型仅数学形式如上。
\[
\boxed{v_1 = \sqrt{\frac{(\rho_2-\rho_1)gl}{\rho_2}}}
\]
\item \textbf{最大浸入深度 $h$（末速度为0，$\rho_2<\rho_1/2$）}
全程重力做功：$W_G = mg \cdot \dfrac{h}{2}$。
浮力分两段：$0\to l$ 变浮力；$l\to h$ 完全浸入，恒浮力 $F_\text{浮}=\rho_1 g S l$。
\[
W_\text{浮} = -\int_{0}^{l}\rho_1 g S x dx - \int_{l}^{h}\rho_1 g S l \,dx
= -\frac12 \rho_1 g S l^2 - \rho_1 g S l (h-l)
\]
动能定理：$W_G + W_\text{浮}=0$
\[
\frac12 \rho_2 S l g h - \rho_1 g S l \left(h-\frac{l}{2}\right)=0
\]
约去公因子：
\[
\frac12 \rho_2 h - \rho_1\left(h-\frac{l}{2}\right)=0
\implies h = \frac{\rho_1 l}{2(\rho_1-\rho_2)}
\]
\[
\boxed{h_\text{max} = \frac{\rho_1 l}{2(\rho_1-\rho_2)}}
\]
\item \textbf{最大速率（合力为0，加速度 $a=0$）}
设浸入深度 $x_0$ 受力平衡：$mg = \rho_1 g S x_0$
\[
\rho_2 S l g = \rho_1 g S x_0 \implies x_0 = \frac{\rho_2}{\rho_1}l
\]
下落 $x_0$ 过程动能定理：
\[
mg\cdot \frac{x_0}{2} - \int_{0}^{x_0}\rho_1 g S x dx = \frac12 m v_\text{max}^2
\]
代入 $m=\rho_2 S l,\ x_0=\dfrac{\rho_2}{\rho_1}l$，化简：
\[
\boxed{v_\text{max} = \sqrt{\frac{\rho_2 g l^2}{\rho_1}\left(2-\frac{\rho_2}{\rho_1}\right)}}
\]
\end{enumerate}
\end{solution}
\subsection{F‑t图像 动量定理求7s末速度、坐标}
一质点沿 $x$ 轴运动，其所受的力如题图所示，设 $t=0$ 时，$v_0=5\ \mathrm{m\cdot s^{-1}}$，$x_0=2\ \mathrm{m}$，质点质量 $m=1\ \mathrm{kg}$，试求该质点 $7\ \mathrm{s}$ 末的速度和坐标。
\begin{solution}
动量定理：$\displaystyle \int_{0}^{t} F\,dt = m(v-v_0)$，$F$‑$t$ 图像面积为冲量。
分段力函数：
\[
\begin{cases}
0\le t\le 5\colon & F_1 = 2t \\
5< t\le 7\colon & F_2 = 35 -5t
\end{cases}
\]
\textbf{0~5 s：}
\[
I_1=\int_0^5 2t\,dt = \big[t^2\big]_0^5=25\ \mathrm{N\cdot s}
\]
\[
a_1=\frac{F_1}{m}=2t,\quad
v_1(t) = v_0+\int_0^t 2t'dt' = 5 + t^2
\]
\[
x_1(t)=x_0+\int_0^t v_1(t')dt' = 2 +5t+\frac{t^3}{3}
\]
$t=5\ \mathrm{s}$：
\[
v_5=5+5^2=30\ \mathrm{m\cdot s^{-1}},\quad
x_5=2+25+\frac{125}{3}=\frac{206}{3}\ \mathrm{m}
\]
\textbf{5~7 s：}
\[
I_2=\int_{5}^{7}(35-5t')dt'
=\left[35t'-\frac52 t'^2\right]_{5}^{7}
=(245-122.5)-(175-62.5)=10\ \mathrm{N\cdot s}
\]
总冲量 $I=I_1+I_2=25+10=35\ \mathrm{N\cdot s}$。
\[
v_7 = v_0+\frac{I}{m}=5+35=40\ \mathrm{m\cdot s^{-1}}
\]
速度函数 $t\in[5,7]$：
\[
v_2(t)=v_5+\int_{5}^{t}(35-5t')dt'
=30+\left(35(t-5)-\frac52(t^2-25)\right)
\]
位置：
\[
x_7 = x_5+\int_{5}^{7} v_2(t')dt' =\frac{418}{3}\approx 139.33\ \mathrm{m}
\]
\[
\boxed{v_7=40\ \mathrm{m\cdot s^{-1}},\quad x_7=\frac{418}{3}\approx 139.33\ \mathrm{m}}
\]
\end{solution}
\subsection{连接体，绳子绷紧后C开始运动（动量守恒）}
三物体质量都为 $m$，B 和 C 紧靠在一起放在光滑的水平面上，二者间连有一段长为 $0.4\ \mathrm{m}$ 的绳子，B 的另一侧的绳子通过滑轮和物体 A 相连，若滑轮和绳子质量不计，滑轮轴上摩擦也不计，绳长一定，问 A，B 开始运动后，经多少时间 C 也开始运动？运动速度为多少？
\begin{solution}
第一阶段：A、B 一起匀加速，C 静止，位移 $s=0.4\ \mathrm{m}$。
对 A、B 整体：$mg = 2ma \implies a=\dfrac{g}{2}$。
匀加速 $s=\dfrac12 a t^2$：
\[
0.4 = \frac12 \cdot \frac{g}{2} \cdot t^2
\implies t=\sqrt{\frac{1.6}{g}}
\]
取 $g=9.8\ \mathrm{m\cdot s^{-2}}$：
\[
\boxed{t \approx 0.404\ \mathrm{s}}
\]
绳子绷紧前瞬间 A、B 速度：
\[
v_0 = a t = \sqrt{2 a s} = \sqrt{2\cdot \frac{g}{2}\cdot 0.4}=\sqrt{0.4g}
\]
绷紧瞬间内力远大于外力，ABC 系统水平动量守恒：
\[
2m v_0 = 3m v
\implies v=\frac23 v_0
\]
\[
\boxed{v = \frac{2}{3}\sqrt{0.4g}\approx 1.32\ \mathrm{m\cdot s^{-1}}}
\]
\end{solution}
\subsection{浮吊水平移动（水平方向质心守恒）}
浮吊，质量为 $m'=2\times 10^4\ \mathrm{kg}$，由岸上吊起质量 $m=2\times 10^3\ \mathrm{kg}$ 的重物后，再将吊杆与竖直方向的夹角 $\theta$ 由 $60^\circ$ 转到 $30^\circ$。设杆长 $l=8\ \mathrm{m}$，水的阻力和杆重忽略不计。求浮吊在水平方向移动的距离。
\begin{solution}
系统水平方向不受外力，水平质心坐标保持不变。
设浮吊相对岸向岸边移动距离为 $s$（浮吊对地水平位移 $-s$）；
重物相对浮吊的水平位移大小：$l(\sin60^\circ-\sin30^\circ)$，
重物相对地面水平位移：$l(\sin60^\circ-\sin30^\circ)-s$。
质心守恒：
\[
m'(-s)+m\big[l(\sin60^\circ-\sin30^\circ)-s\big]=0
\]
\[
s(m'+m)=ml\left(\frac{\sqrt{3}}{2}-\frac12\right)
\]
\[
s = \frac{ml(\sqrt{3}-1)}{2(m'+m)}
\]
代入数值：
\[
s=\frac{2\times10^3 \times 8 (\sqrt{3}-1)}{2(2\times10^4+2\times10^3)}
=\frac{8(\sqrt{3}-1)}{22}\approx 0.266\ \mathrm{m}
\]
\[
\boxed{s\approx 0.266\ \mathrm{m}，浮吊向岸边移动约0.266\ \mathrm{m}}
\]
\end{solution}
